\section{Locating physical choices in the existing macro architecture}
\label{sec:architecture-placement}
MiniLinq already exposes different units of replacement: a CQ planning pass,
a named RA program, an indexed physical lowering, a local Access quotation,
and a Pipeline realization. The concrete hash backend is one implementation
of those boundaries. Its data structures do not define the extensibility
limit of the architecture. The design in the following section is one
additional backend family, organized around immutable row arenas. Its proposed
overlay is not an approved IR addition; necessity and stage placement remain obligations.

\subsection{What the implementation actually exposes}
\code{crates/mini-linq/src/lib.rs} provides these direct entries;
\code{cq-compiler} retains flat typed pass exports. Its \code{physical} module
contains sibling materialized, lazy and composed implementations; \code{pull.rs}
and \code{pipeline.rs} own the lower stages.
The ordinary Rust AST returned by a physical pass owns allocation, construction,
execution order and residual macro calls. Access contains native Rust sources,
predicates and output expressions rather than a fixed collection vocabulary.
Its label states logical binding information; it does not secretly implement
an index. The same source-level Access pattern can consequently be instantiated
with distinct lawful Rust expressions without a shared data-structure trait.

\begin{center}\small
\begin{tabular}{>{\raggedright\arraybackslash}p{.28\linewidth}>{\raggedright\arraybackslash}p{.31\linewidth}>{\raggedright\arraybackslash}p{.31\linewidth}}
\toprule
Existing entry / boundary & Unit the author controls & Appropriate physical variation\\
\midrule
\code{mini_linq!}, \code{planned_query!} / CQ passes & Planning a whole query before committing to named RA & Choose a legal logical schedule or derive a checked YA region.\\
\code{relational_plan!}, \code{planned_relational_plan!} / RA passes & Independently authored named algebra & Choose exact joins, reductions, projections and shared definitions.\\
\code{index_requirements!}, \code{materialized!} / physical passes & Rust file containing storage and local Access holes & Change build/probe sequencing, view storage, runtime specialization, or realize an entire certified region.\\
\code{pull!} / Access & Explicit Rust sources and predicates with logical binding labels & Substitute exact membership expressions; consume a filtered or block-aware source whose construction occurs outside the hole.\\
\code{iterator_pipeline!} / Pipeline & Graph of Scan, Join, Filter, Project and Distinct & Specify an existing flat pipeline directly, or replace its realization pass for an admitted domain.\\
\bottomrule
\end{tabular}
\end{center}
Entries parse their source directly and call the next macro through a relative
\code{mini_linq} path. A caller-owned module facade aliased as
\code{mini_linq} can bind these continuations through ordinary re-exports;
without that alias, the crate supplies the defaults.
\code{materialized!} defaults its local Access holes to relative
\code{mini_linq::collect_access} and accepts an explicit \code{access path;} override;
its checker validates the configured path with the storage and clause contracts.
The composed RA entry now preserves the source in an empty IndexRequirements
overlay, forwarding it to relative \code{mini_linq::composed!}. That typed edge
rejects written requirements and leaves auxiliary construction to recipes.
Compact contextual inputs refine realization; public typed boundaries return
diagnostics for invalid source IR in both debug and release builds.

\paragraph{Ordinary aliases bind compatible implementations.}
The facade must be in the macro invocation's scope; importing just the entry
macro under a new name does not rebind the generated continuation. Nested modules
import the facade explicitly. Absolute \code{::mini_linq} paths in the facade
and recipes reach concrete implementations without recursing through overrides.
The caller may also supply a stable imported name, for example
\code{use ::mini_linq::collect_access as access;}, as the Access path. An alternative macro
can occupy that alias when it accepts the same source language and establishes
the boundary's result, evaluation and ownership contract. Generated paths must
actually reference the alias; an unrelated local import cannot redirect a
literal \code{::mini_linq::pull!} call, as retained in the legacy basic/lazy
physical emitters. The materialized default instead uses the set endpoint.
Configurable physical entries can instead supply
\code{access mini_linq::collect_access;}. An alias at an established language
boundary requires no additional syntax. Whole-operator replacement instead needs
the retained RA definition and its physical binding assumptions: the next section
declares the retained definition with only its free bindings and local laws.
Key construction uses existing Project; reduction uses Semijoin on that projected
resource. Configured lexical macros bind storage, continuation and optional work
once, outside the semantic inputs. Operator control precedes Access
lowering, while provider calls in native leaves can expand later. The set endpoint
folds terminal Distinct into its collection sink; maps select providers and recipes.
Compatibility still depends on the stated semantic and ownership contracts.
Within composed compilation, layout facts, resource planning, query API construction
and invocation auditing occupy separate modules. Provider macros own each binding's
type and access; operator macros own construction and traversal. Hash and nested-loop
joins are selectable without changing semijoin or API lowering. Native sets coexist
with block storage. The data structures are in \code{cq_storage}; generated code
associates storage and masks without a universal runtime relation interface.
Unobserved execution omits statistics operations; the observed entry shares the
same lowering with counters enabled.

\subsection{Reading the papers at the right granularity}
The literature separates representation, scheduling and intermediate semantics.
The right column below is our architectural placement judgment, not a claim
that those systems use MiniLinq's macros.
\begin{center}\small
\begin{tabular}{>{\raggedright\arraybackslash}p{.24\linewidth}>{\raggedright\arraybackslash}p{.35\linewidth}>{\raggedright\arraybackslash}p{.31\linewidth}}
\toprule
Primary source & Mechanism reported & Placement in MiniLinq\\
\midrule
SQL Server, CIDR 2026, \S\S3.1,4.1~\cite{ya-sqlserver} & Optimizer decides whether filtering is worthwhile; hash-build execution chooses Bloom/exact bit vector and size from observed keys. & Keep static eligibility in the pass; emit runtime construction and branch selection in its Rust context.\\
Predicate Transfer, CIDR 2024, \S3.2~\cite{ya-pt} & Prototype fixes Bloom representation and propagates filters between relations, rebuilding them when key columns change. & A complete transfer region needs its own schedule and candidate invariant; reuse local Access for legal row operations.\\
Diamond Hardened Joins, PVLDB 2024, \S3.2~\cite{ya-dhj} & Separates exact Lookup from Expand, retaining handles to matching tuples so their enumeration can occur later. & Change operator sequencing and handle lifetime in the physical region; a membership-expression replacement alone does not implement it.\\
YA in Column Stores, PVLDB 2025, \S\S3--4~\cite{ya-shredded} & Gives nested operators and a shredded realization with grouped references, while reusing applicable flat operators. & A backend-specific nested language may be useful; preserve compatible lower holes without forcing nested results into the current flat row-set contract.\\
\bottomrule
\end{tabular}
\end{center}
SQL Server's instance-optimality discussion assumes exact bitmap filtering;
its practical Bloom choice does not automatically carry that bound. Likewise,
nested/bag results do not inherit MiniLinq's set-projection laws. These are
boundaries of the associated proofs, even when implementations share a method
name or all return the same final query answer.

\subsection{Replacing an expression, a context, or a compiler segment}
For a represented exact key set $S$, a hash set and an injectively encoded
bit vector can both provide an expression $e(k)$ satisfying
$e(k)\iff k\in S$. This licenses substitution at a membership leaf when
borrows, binding types, effects and termination match its surrounding context.
A Bloom expression satisfies only $k\in S\Rightarrow e(k)$, so it cannot
replace that exact leaf. It can precede an exact test, or participate in a
whole-region conservative transfer whose final assembly is exact.

The manuscript's ST-Hole rule keeps the surrounding Rust context fixed.
Changing storage ownership, scheduling, publication barriers or access granularity
therefore requires replacing the owning context as well. For example, a
zone-map source must decide block rejection before reading the skipped payload.
Installing a clever \code{for} source after the existing eager input
\code{HashSet} normalization cannot save that earlier input read. A block-aware
physical pass must emit the appropriate input/scan context; it may still reuse
\code{pull!} for the local output loop.

For a strategy $a$, write its compiler edge as
\[
 f_a:L\longrightarrow\mathsf{Rust}[L_1,\ldots,L_m],\qquad
 f_a(s)=C_a[h_1,\ldots,h_m].
\]
Here $L$ is an existing input language, each $h_j$ an independently checked
lower-language hole, and $C_a$ the generated runtime context. The strategy
declares its representation relation $\gamma_a$ and admitted domain; it must
preserve the observation $O$ and client class required by the replaced boundary.
The manuscript's conditional composition theorem applies with all its premises:
total compilation on the declared domain, valid targets admitted by subsequent
passes, established runtime preconditions, compatible environment relations,
and adequate closing contexts. There need be
neither one common physical intermediate representation for all $a$ nor one
trait shared by their storage objects. When two strategies need different
intermediate meanings, their final observations can still agree through their
separate representation relations and whole-region proofs.

This yields a concrete implementation discipline: choose the smallest existing
boundary that owns all changed semantics. Use a native expression for a lawful
local replacement; replace the Rust context when construction or lifetime
changes; replace a larger pure RA region when reduction or intermediate-result
semantics changes. A new language needs a demonstrated missing distinction,
formal semantics, local contracts and pass rules before IR implementation.
Its position must respect CQ, RA, IndexRequirements, the Rust context with Access,
and Pipeline in that order; private token protocols do not satisfy this discipline.

\subsection{Static expansion can emit a runtime representation choice}
Consider an exact semijoin whose build keys may occupy a compact integer range.
A physical pass can emit two ordinary Rust branches, each containing an Access
quotation with the same logical bindings but different concrete expressions:
\begin{lstlisting}
let answer = if observed_range_fits_bitmap(&build_keys) {
    let bits = build_exact_bitmap(&build_keys);
    mini_linq::pull! {
        out(k) => ((k,)) :-
        Probe(k) => for (k,) in probe_rows.iter().copied(),
        Build(k) => if bitmap_contains(bits, k);
    }.collect::<HashSet<_>>()
} else {
    let keys = build_exact_hash_set(&build_keys);
    mini_linq::pull! {
        out(k) => ((k,)) :-
        Probe(k) => for (k,) in probe_rows.iter().copied(),
        Build(k) => if keys.contains(&k);
    }.collect::<HashSet<_>>()
};
\end{lstlisting}
The compiled witness is \path{example/src/bin/access_representations.rs};
this listing abbreviates its helpers. Both macro invocations expand
at compile time. The generated Rust chooses one branch from input data during
execution. Branch-local storage types need not unify: each branch returns the
same declared answer type. A trait, enum or shared helper remains an optional
implementation choice when it would reduce code or provide another useful
interface. The witness's range cutoff is illustrative, not a measured optimizer.

For correctness, let $E_b,E_h$ be the completed bitmap and hash branches,
and $\chi(I)$ the runtime selector. Require
\[
\begin{aligned}
 \chi(I)&\Rightarrow \sem{E_b}_I=\sem{l\ltimes r}_I,\\
 \neg\chi(I)&\Rightarrow \sem{E_h}_I=\sem{l\ltimes r}_I.
\end{aligned}
\]
The conditional then has that denotation regardless of the selector's
profitability. The selector must be total and meet the host assumptions;
its successful branch test must establish representation-specific conditions
such as bounded arithmetic and injective bitmap addressing. This is ordinary
conditional refinement and contextual adequacy, not a new meaning for RA.
At the weaker filter interface, exact-bit-vector and Bloom branches can instead
both establish Cover, provided the caller still performs the exact operation
required by its own contract.

\paragraph{What the following physical design promises.}
PA-Choose selects compatible resources within the stated arena backend family.
Its static \code{Provider} enumeration is private planning data, not a required
runtime trait or a universal catalog for every compiler entry. The explicit
Exact/Transfer modes specify two different observation contracts. Factorized
YA, persistent block storage and additional execution strategies may
replace an appropriate compiler segment and reuse compatible existing holes;
none must first be coerced into the arena representation.
The current default physical checker deliberately recognizes the default
emitter's templates. An alternative edge needs its own source-to-target
contract and checker. Macro checks cover local structure and bindings; Rust
checks the completed expansion's types. Neither proves arbitrary provider laws.
The executable witness checks selected
branch laws and existing macro composition, not arbitrary contextual adequacy.
